\section{Experimental Design}
\label{sec:experiments}

We design five controlled experiments (A--E) to systematically characterize the Z-layer selection space, each isolating a single variable while holding all others fixed. All experiments use Scroll~1 (PHerc.Paris.4) segments (65 TIF layers, 7.91~\textmu m resolution) with a common base configuration: tile size $P = 64$, training stride $\delta = 32$, batch size $B = 32$, learning rate $\eta = 3 \times 10^{-4}$, warmup-cosine schedule (3-epoch warmup), intensity clip $\tau = 200$, and evaluation threshold $\theta = 0.4$. The test segment (20230504223647) is held out and never used during training or hyperparameter selection. Ground truth is provided by pseudo-labels generated from the GP TimeSformer.

\subsection{Experiment A: Layer-Count Ablation}
\label{sec:exp_a}

\textbf{Question:} At what point does adding more Z-layers stop improving ink detection?

A fresh UNet2D (no ImageNet pre-training) is trained for 30 epochs for each $N \in \{3, 5, 7, 9\}$, with reference to Table~\ref{tab:ablation_configs} the window centered symmetrically on layer~32 using $s = 32 - \lfloor N/2 \rfloor$. All parameters except $N$ (and $s$) are held constant; the 100-epoch ImageNet-pretrained baseline at $N = 5$ serves as reference.

\begin{table}[h]
\centering
\caption{Experiment A: Layer-count ablation configurations}
\label{tab:ablation_configs}
\begin{tabular}{cccc}
\hline
$N$ & Start Layer & End Layer & Physical Span (\textmu m) \\
\hline
3  & 31 & 33 & 23.7 \\
5  & 30 & 34 & 39.6 \\
7  & 29 & 35 & 55.4 \\
9  & 28 & 36 & 71.2 \\
\hline
\end{tabular}
\end{table}

\subsection{Experiment B: Depth Window Position Sweep}
\label{sec:exp_b}

\textbf{Question:} Which depth window within the 65-layer stack produces the strongest ink signal?

With $N = 5$ fixed, the start layer is swept across $s \in \{28, 29, 30, 31, 32\}$ with reference to Table~\ref{tab:position_sweep} using the pre-trained 100-epoch UNet2D in inference-only mode (no retraining per position). This isolates the effect of depth placement from training dynamics; all model weights and inference parameters are identical across configurations.

\begin{table}[h]
\centering
\caption{Experiment B: Position sweep configurations}
\label{tab:position_sweep}
\begin{tabular}{ccl}
\hline
Start $s$ & Layers & Anatomical Region \\
\hline
28 & 28--32 & Surface approach zone \\
29 & 29--33 & Crossing into ink zone \\
\textbf{30} & \textbf{30--34} & \textbf{Core ink zone (expected optimum)} \\
31 & 31--35 & Surface center to sub-surface \\
32 & 32--36 & Below surface center \\
\hline
\end{tabular}
\end{table}

\subsection{Experiment C: Leave-One-Out Layer Sensitivity}
\label{sec:exp_c}

\textbf{Question:} Which individual layers within the 5-layer window contribute most to ink detection?

Each of the five input layers (30-34) is individually zeroed, and the mean absolute prediction change is measured. Model weights are unchanged; only the input data varies. Results are reported over three segments (test, validation, one training) to assess cross-segment consistency.

\subsection{Experiment D: Cross-Segment Generalization}
\label{sec:exp_d}

\textbf{Question:} Does the default configuration ($N = 5$, $s = 30$) generalize without per-segment tuning?

The 100-epoch UNet2D is applied identically to all 8 ground-truth segments of Scroll~1 (4 training, 1 validation, 1 test, 2 additional GT segments). Cross-segment stability is quantified by the spread of $F_{0.5}$ scores across these segments.

\subsection{Experiment E: Failure Analysis}
\label{sec:exp_e}

\textbf{Question:} What causes the default configuration to fail on specific segments?

The worst-performing segments from Experiment~D are diagnosed via depth profile inspection, coverage statistics, and prediction standard deviation checks. Each failure is classified into one of three hypothesized categories: depth offset (ink signal outside the 30-34 window), spatial misalignment (predictions present but not overlapping ground truth), or low contrast (insufficient ink signal).

\subsection{Benchmark Comparison}
\label{sec:benchmark}

In addition to the five experiments, our UNet2D and TimeSformer models (trained for 30, 80, and 100 epochs) are directly compared against the GP Winner on the held-out test segment, establishing whether a focused 5-layer approach can match or surpass the original 26-layer GP model and at what computational cost.
\vspace{2\baselineskip}

\newpage
